% !TEX program = pdflatex
% Odpowiedzi na pytania szczegółowe do prezentacji Urban Health ABM.
% Kompilacja: pdflatex NOTATKI_QA.tex (×2 dla spisu treści).

\documentclass[11pt,a4paper]{article}

% ===== JĘZYK =====
\usepackage[T1]{fontenc}
\usepackage[utf8]{inputenc}
\usepackage[polish]{babel}
\usepackage{lmodern}
\usepackage{microtype}

% ===== MATEMATYKA =====
\usepackage{amsmath, amssymb, amsthm, mathtools}
\usepackage{bm}

% ===== TABELE =====
\usepackage{booktabs}
\usepackage{array}
\usepackage{tabularx}

% ===== LAYOUT =====
\usepackage[a4paper, margin=2.5cm, headheight=22pt]{geometry}
\usepackage[dvipsnames, table]{xcolor}
\usepackage{wasysym}
\usepackage{fancyhdr}
\usepackage{titlesec}
\usepackage{enumitem}
\usepackage{tcolorbox}
\tcbuselibrary{breakable, skins}
\usepackage{listings}
\usepackage[hidelinks]{hyperref}
\hypersetup{colorlinks=true, linkcolor=NavyBlue!70!black, urlcolor=NavyBlue!70!black}

% ===== KOLORY (spójne z poprzednim plikiem) =====
\definecolor{ink}{HTML}{1B2A41}
\definecolor{gold}{HTML}{A07840}
\definecolor{teal}{HTML}{2C5F7C}
\definecolor{cream}{HTML}{F6F1E8}
\definecolor{rule}{HTML}{D8DEE5}
\definecolor{muted}{HTML}{6B7785}
\definecolor{codebg}{HTML}{F4F2EC}
\definecolor{warmred}{HTML}{C0392B}

% ===== TYTUŁY =====
\titleformat{\section}[hang]
  {\Large\bfseries\color{ink}}
  {\thesection}{0.7em}{}[\vspace{2pt}\textcolor{gold}{\hrule height 1.5pt}\vspace{2pt}]
\titleformat{\subsection}[hang]
  {\large\bfseries\color{ink}}
  {\thesubsection}{0.6em}{}
\titleformat{\subsubsection}[hang]
  {\bfseries\color{teal}}
  {\thesubsubsection}{0.5em}{}

% ===== GŁÓWKA =====
\pagestyle{fancy}
\fancyhf{}
\fancyhead[L]{\small\color{muted}\textsc{Urban Health ABM --- Pytania i odpowiedzi}}
\fancyhead[R]{\small\color{muted}\thepage}
\renewcommand{\headrulewidth}{0.4pt}
\renewcommand{\headrule}{\hbox to\headwidth{\color{rule}\leaders\hrule height \headrulewidth\hfill}}

% ===== KOD =====
\lstdefinestyle{py}{
  language=Python,
  basicstyle=\small\ttfamily\color{ink},
  backgroundcolor=\color{codebg},
  keywordstyle=\color{teal}\bfseries,
  stringstyle=\color{gold!90!black},
  commentstyle=\color{muted}\itshape,
  numberstyle=\tiny\color{muted},
  showstringspaces=false,
  frame=leftline,
  framerule=1.5pt,
  rulecolor=\color{teal},
  xleftmargin=12pt,
  framexleftmargin=10pt,
  breaklines=true,
  numbers=left,
  numbersep=8pt,
  columns=fullflexible,
  belowskip=10pt,
  aboveskip=10pt
}
\lstset{style=py}

% ===== BOXY =====
\newtcolorbox{infobox}[1][]{
  colback=cream, colframe=gold, boxrule=0pt, leftrule=2pt, arc=2pt,
  left=10pt, right=10pt, top=6pt, bottom=6pt, fontupper=\small, breakable, #1
}
\newtcolorbox{warnbox}[1][]{
  colback=warmred!4!white, colframe=warmred, boxrule=0pt, leftrule=2pt, arc=2pt,
  left=10pt, right=10pt, top=6pt, bottom=6pt, fontupper=\small, breakable, #1
}
\newtcolorbox{formulabox}[1][]{
  colback=white, colframe=rule, boxrule=0.5pt, arc=2pt,
  left=10pt, right=10pt, top=6pt, bottom=6pt, fontupper=\normalsize, breakable, #1
}
\newtcolorbox{qbox}[1][]{
  colback=teal!8!white, colframe=teal, boxrule=0pt, leftrule=3pt, arc=2pt,
  left=12pt, right=10pt, top=8pt, bottom=8pt,
  fontupper=\normalsize\bfseries\color{teal}, breakable, #1
}

% ===== KOMENDY =====
\newcommand{\promile}{\ensuremath{\text{\textperthousand}}}

% ===== METADANE =====
\title{
  \vspace{-1.5cm}
  {\small\color{gold}\textsc{Notatki uzupełniające --- projekt zespołowy 2026}}\\[10pt]
  {\Huge\bfseries\color{ink} Pytania i odpowiedzi}\\[6pt]
  {\large\color{ink!70} Urban Health ABM --- detale parametrów, walidacji i implementacji}\\[20pt]
  {\color{rule}\hrule height 1pt}
}
\author{}
\date{}

\begin{document}
\maketitle
\thispagestyle{empty}
\vspace{1cm}

\noindent\textcolor{muted}{\small Dokument towarzyszący prezentacji. Odpowiada na pięć szczegółowych pytań dotyczących:
parametrów FM i MM (zakresy, jednostki, wyprowadzenie),
zakresu suwaków czynników ryzyka [0,\,3],
przykładu obliczeniowego z 60-latkiem,
rozbieżności CBR modelu vs GUS,
oraz struktur danych Citizen i Zone.}

\newpage
\tableofcontents
\newpage

% ============================================================
\section{Jak jest wyliczane FM i MM?}
% ============================================================

\begin{qbox}
Co oznaczają FM i MM, skąd się biorą, dlaczego badane są w takich przedziałach? Jakie konkretne wartości testujemy?
\end{qbox}

\subsection{Definicje}

\textbf{FM} (\emph{Fertility Multiplier}) oraz \textbf{MM} (\emph{Mortality Multiplier}) to \textbf{bezwymiarowe mnożniki} skalujące całe tabele demograficzne z GUS Polska 2021. Nie są stopami w promilach --- to mnożniki ($\times$).

\subsubsection{FM --- mnożnik płodności}

\begin{formulabox}
\[
P_{\text{birth, monthly}}(\text{age}) = \frac{\text{ASFR}(\text{age})}{12} \cdot \text{FM} \cdot \text{disease\_reduction}
\]
\end{formulabox}

W kodzie (\texttt{simulation\_engine.py}):
\begin{lstlisting}
monthly_birth_prob = ASFR[wiek_kobiety] / 12 * fertility_multiplier
\end{lstlisting}

\textbf{FM = 1.0} odpowiada bezpośrednio polskim danym GUS 2021 (TFR\,=\,1.28, CBR\,=\,8.30\promile). Mnożnik różny od 1 skaluje \emph{każdą} z 7 wiekowych stóp ASFR (15--19, 20--24, \ldots, 45--49 lat) tym samym współczynnikiem.

\subsubsection{MM --- mnożnik śmiertelności}

\begin{formulabox}
\[
P_{\text{death, monthly}} = q_x(\text{age, sex}) \cdot \text{disease\_mult} \cdot \text{cox\_mult} \cdot \text{MM}
\]
\end{formulabox}

W kodzie:
\begin{lstlisting}
monthly_death_prob = q_x[wiek][plec] * disease_mult * cox_mult * mortality_multiplier
\end{lstlisting}

\textbf{MM = 1.0} odpowiada Polsce 2021 (CDR\,=\,15.6\promile). MM skaluje każdą $q_x$ z tabeli śmiertelności GUS dla każdego przedziału wiekowego i płci.

\subsection{Skąd się biorą --- skala bazowa}

Bazowe wartości pochodzą bezpośrednio z tabel GUS:

\begin{itemize}[leftmargin=*]
  \item \textbf{ASFR} --- 7 wiekowych stóp płodności (15--19, 20--24, \ldots, 45--49 lat); razem dają TFR\,=\,1.28.
  \item \textbf{$q_x$} --- 10 wiekowych prawdopodobieństw zgonu × 2 płcie z Tablic Trwania Życia 2021.
\end{itemize}

Te bazowe stopy są dane raz na zawsze. \textbf{FM i MM są jedynymi zmiennymi globalnymi}, które pozwalają eksperymentować z różnymi scenariuszami demograficznymi.

\subsection{Dlaczego mnożnik, a nie stopa?}

W pierwszej wersji modelu parametry nazywały się \texttt{fertility\_rate} i \texttt{mortality\_rate}. Były mylące --- sugerowały jednostkę "stopa" (\promile{}). W rzeczywistości to bezwymiarowe mnożniki tabel wiekowych.

\textbf{Stara nazwa $\to$ pomyłki:}
\begin{itemize}
  \item \texttt{fertility\_rate = 2.0} $\to$ ktoś czytał "TFR = 2.0" lub "CBR = 2\promile{}". Naprawdę: TFR $\approx$ 2.56, CBR $\approx$ 16.6\promile{}.
  \item \texttt{mortality\_rate = 0.5} $\to$ wygląda jak "0.5/1000/rok". Naprawdę: "połowa polskiej śmiertelności w każdym wieku".
\end{itemize}

Po refaktorze:
\begin{itemize}
  \item \textbf{mnożnik} --- jednostka ($\times$), bezwymiarowy
  \item \textbf{stopa} --- jednostka (\promile{}), opisuje konkretną wartość demograficzną
\end{itemize}

\subsection{Zakres siatki w gridsearch}

\subsubsection{FM $\in [0.40,\ 2.50]$ z 12 punktami}

\begin{formulabox}
\[
\text{FM} \in \texttt{np.linspace(0.4, 2.5, 12)} = \{0.40,\ 0.59,\ 0.78,\ 0.97,\ 1.16,\ 1.35,\]
\[1.54,\ 1.73,\ 1.92,\ 2.12,\ 2.31,\ 2.50\}
\]
\end{formulabox}

\textbf{Dlaczego ten zakres?}
\begin{itemize}[leftmargin=*]
  \item \textbf{Dolna granica 0.40} --- daje TFR $\approx$ 0.51, czyli skrajny scenariusz depopulacji.
  \item \textbf{Górna granica 2.50} --- daje TFR $\approx$ 3.20, przekracza wyraźnie próg zastępowalności (2.10), pozwala badać scenariusze silnego wzrostu.
  \item \textbf{Stabilność leży wewnątrz} --- z reguły proxy $\text{FM}_{\text{stable}} = 1.88 \cdot \text{MM}$, więc dla $\text{MM} \in [0.3, 1.6]$ stabilne FM mieści się w \mbox{[0.56, 3.0]}. Wybrany zakres pokrywa większość krzywej stabilności.
  \item \textbf{12 punktów} --- balans między rozdzielczością a kosztem obliczeniowym ($12^2 = 144$ symulacji $\times$ 90\,s każda).
\end{itemize}

\subsubsection{MM $\in [0.30,\ 1.60]$ z 12 punktami}

\begin{formulabox}
\[
\text{MM} \in \texttt{np.linspace(0.3, 1.6, 12)} = \{0.300,\ 0.418,\ 0.536,\ 0.655,\ 0.773,\]
\[0.891,\ 1.009,\ 1.127,\ 1.245,\ 1.364,\ 1.482,\ 1.600\}
\]
\end{formulabox}

\textbf{Dlaczego węższy zakres niż FM?}
\begin{itemize}[leftmargin=*]
  \item \textbf{Dolna granica 0.30} --- skrajna redukcja śmiertelności (hipotetyczny przełom medyczny).
  \item \textbf{Górna granica 1.60} --- 60\,\% wzrost śmiertelności.
  \item \textbf{Asymetria względem FM} --- stabilność wymaga FM $>$ MM (proxy: FM $\approx$ 1.88\,$\cdot$\,MM). Przy MM = 1.6 stabilne FM byłoby już $\approx$ 3.0 (górna granica siatki FM), więc dalsze rozszerzanie MM nie wniosłoby informacji.
\end{itemize}

\subsection{Interpretacja kluczowych wartości}

\begin{table}[h]
\centering
\small
\begin{tabular}{rrlrr}
\toprule
\textbf{FM} & \textbf{Co dzieje się z ASFR} & & \textbf{TFR} & \textbf{CBR} \\
\midrule
0.5 & każda $\div$ 2     & (kampania anty-natalistyczna) & 0.64 & 4.2\promile \\
1.0 & bez zmian (GUS)    & polska baseline               & 1.28 & 8.30\promile \\
1.88 & każda $\times$ 1.88 & punkt stabilności proxy      & 2.41 & 15.6\promile \\
2.0 & każda $\times$ 2   & podwojenie                    & 2.56 & 16.6\promile \\
2.5 & każda $\times$ 2.5 & granica górna siatki          & 3.20 & 20.8\promile \\
\midrule
\textbf{MM} & \textbf{Co dzieje się z $q_x$} & & \textbf{CDR} & \textbf{Śr. dł. życia} \\
\midrule
0.3 & każda $\times$ 0.3 & przełom medyczny              & 4.7\promile & bardzo długie \\
1.0 & bez zmian (GUS)    & polska baseline               & 15.6\promile & $\sim$77 lat \\
1.6 & każda $\times$ 1.6 & granica górna siatki          & 25\promile & znacznie krótsze \\
\bottomrule
\end{tabular}
\caption{Interpretacja wybranych wartości FM i MM.}
\end{table}

\subsection{Punkt stabilności znaleziony przez pełny ABM}

\begin{infobox}
\textbf{Wynik:} po przeszukaniu siatki 144 punktów pełnym ABM (50\,k agentów $\times$ 50 lat każdy), punkt stabilności populacji Polski to:
\[
\text{FM} \approx 1.93, \quad \text{MM} \approx 1.01 \quad \implies \quad \text{score} = -0.88\,\%
\]
Czyli Polska potrzebuje \emph{prawie dwukrotnie wyższej płodności} niż obecna (TFR\,$\approx$\,2.47), aby populacja przestała się kurczyć przy obecnej śmiertelności.
\end{infobox}

\newpage
% ============================================================
\section{Dlaczego czynniki ryzyka są w przedziale [0, 3]?}
% ============================================================

\begin{qbox}
Co oznacza \texttt{smoking = 0.3}? Co oznacza \texttt{smoking = 2.7}? Dlaczego granica górna to 3, a nie np.\ 5?
\end{qbox}

\subsection{Suwaki RF w aplikacji Streamlit}

W aplikacji \texttt{interactive\_simulation\_app.py} każdy z 7 czynników ryzyka ma suwak o zakresie $[0.0,\ 3.0]$ z krokiem 0.1. Suwaki to \textbf{mnożniki prevalencji} --- skalują bazowe prawdopodobieństwo wystąpienia danego RF u agenta przy tworzeniu populacji.

\begin{lstlisting}
rf_multipliers[rf] = st.slider(
    f"{rf.replace('_', ' ').title()}",
    min_value=0.0,
    max_value=3.0,
    value=1.0,
    step=0.1,
)
\end{lstlisting}

\subsection{Jak mnożnik działa --- przykład dla palenia}

Bazowe prawdopodobieństwo palenia zależy od wieku agenta (kalibrowane do danych Białystoku+):

\begin{formulabox}
\[
P_{\text{smoking, base}}(\text{age}) = \max\left( 0.25 \cdot \left(1 - \frac{(\text{age} - 45)^2}{50^2}\right),\ 0.10 \right) \quad \text{dla age} \in [20, 70]
\]
\end{formulabox}

Mnożnik suwaka jest stosowany bezpośrednio:
\begin{formulabox}
\[
P_{\text{smoking, effective}}(\text{age}) = \min\left( P_{\text{smoking, base}}(\text{age}) \cdot \text{slider},\ 1.0 \right)
\]
\end{formulabox}

\subsection{Konkretne wartości suwaków}

Dla agenta w wieku szczytowej prevalencji palenia (age = 45, $P_{\text{base}}$ = 0.25):

\begin{table}[h]
\centering
\small
\begin{tabular}{rrrl}
\toprule
\textbf{Suwak} & \textbf{Mnożnik} & \textbf{P efektywne (age=45)} & \textbf{Interpretacja} \\
\midrule
0.0 & $\times$ 0    & 0\,\%   & całkowita eliminacja palenia (scenariusz utopijny) \\
0.3 & $\times$ 0.3  & 7.5\,\% & redukcja o 70\,\% (silna kampania anty-nikotynowa) \\
0.5 & $\times$ 0.5  & 12.5\,\% & redukcja o 50\,\% \\
1.0 & $\times$ 1    & 25\,\%  & polska baseline z Białystoku+ \\
1.5 & $\times$ 1.5  & 37.5\,\% & wzrost o 50\,\% \\
2.7 & $\times$ 2.7  & 67.5\,\% & wzrost o 170\,\% (skrajnie wysokie palenie) \\
3.0 & $\times$ 3    & 75\,\%  & potrojenie (górna granica suwaka) \\
\bottomrule
\end{tabular}
\caption{Interpretacja wartości suwaka \texttt{smoking}.}
\end{table}

\subsection{Dlaczego górna granica to 3, a nie więcej?}

\begin{enumerate}[leftmargin=*]
  \item \textbf{Powyżej 3 prawdopodobieństwa się nasycają.} Dla wielu wieków $P_{\text{base}} \cdot 3$ przekracza 1.0 i jest kapowane do 1.0. Dalsze zwiększanie mnożnika nie wprowadza już zmian.
  \item \textbf{Skrajna eksploracja przestrzeni.} 3 $\times$ baseline odpowiada hipotetycznemu społeczeństwu, w którym RF jest stale obecny u $\sim$70\,\% populacji. To wartość ekstrema dla scenariuszy "co najgorsze może się zdarzyć".
  \item \textbf{Dolna granica 0.0} --- pełna eliminacja (utopia zdrowia publicznego), pozwala mierzyć "potencjał interwencji" --- ile populacji uratujemy w pełnym scenariuszu eliminacji.
  \item \textbf{Symetria interpretacyjna} --- $[0,\ 3]$ z baseline 1.0 daje symetryczny zakres: 1/3 do 3 (redukcja vs amplifikacja $\times$ 3 w obu kierunkach).
\end{enumerate}

\subsection{Przykład: kampania anty-nikotynowa}

\begin{infobox}
\textbf{Scenariusz} --- 70\,\% redukcja palenia (suwak \texttt{smoking = 0.3}), pozostałe RF bez zmian.

\smallskip

\textbf{Spodziewane efekty po 50 latach:}
\begin{itemize}
  \item Znaczący spadek prevalencji raka płuc (ponieważ HR\,=\,15 dla palenia $\to$ rak płuc dominuje w determinacji ryzyka LC).
  \item Spadek CVD (palenie ma HR\,=\,2.5 dla CVD).
  \item Wyższa końcowa populacja (mniej zgonów z chorób nikotyno-zależnych).
  \item Spadek mediany $H_{\mathrm{cum}}$ w populacji palaczy poniżej $\sim$0.05.
\end{itemize}
\end{infobox}

\subsection{Skutek na inicjalizację RF}

W \texttt{interactive\_simulation\_app.py} (funkcja \texttt{modified\_init\_rf}) suwak jest aplikowany na każdy z 7 RF analogicznie:

\begin{lstlisting}
smoking_prob *= rf_multipliers.get("smoking", 1.0)
obesity_prob = min(obesity_prob, 0.45) * rf_multipliers.get("obesity", 1.0)
inactivity_prob *= rf_multipliers.get("physical_inactivity", 1.0)
alcohol_prob = base * rf_multipliers.get("alcohol_abuse", 1.0)
cholesterol_prob *= rf_multipliers.get("high_cholesterol", 1.0)
hypertension_prob *= rf_multipliers.get("hypertension_stage0", 1.0)
family_prob = 0.15 * rf_multipliers.get("family_history", 1.0)
\end{lstlisting}

\newpage
% ============================================================
\section{Przykład z 60-letnim mężczyzną --- skąd wartości i dlaczego wiek bazowy 30?}
% ============================================================

\begin{qbox}
Na slajdzie jest przykład obliczeniowy: 60-letni mężczyzna z paleniem, cholesterolem i nadciśnieniem. Skąd biorą się konkretne liczby? Czy wiek bazowy to 30 lat? Dlaczego akurat 30?
\end{qbox}

\subsection{Pełny wzór Coxa}

\begin{formulabox}
\[
\Delta h_d(a, X) = \lambda_{0,d} \cdot \exp\bigl(\gamma_{\text{age},d} \cdot (a - 30)\bigr) \cdot \exp\!\left(\sum_{i=1}^{7} \beta_{d,i} \cdot X_i\right)
\]
\end{formulabox}

\subsection{Wartość po wartości}

Wszystkie liczby pochodzą z macierzy zdefiniowanych w \texttt{disease\_model.py} (linie 46--78):

\subsubsection{$\lambda_{0,\text{CVD}} = 6.0 \times 10^{-5}$ --- bazowy hazard CVD}

Z \texttt{BASELINE\_HAZARD["CVD"]}. Jest to miesięczne ryzyko inicjacji CVD u \emph{30-letniej osoby bez żadnego RF}. Wartość skalibrowana do danych WHO/ESC tak, aby populacyjna prevalencja CVD zgadzała się z polską ($\sim$35\,\% dorosłej populacji).

\subsubsection{Współczynnik wieku: $\exp(0.06 \cdot 30) \approx 6.05$}

Z \texttt{AGE\_HAZARD\_GROWTH["CVD"] = 0.06} oraz $a - 30 = 30$ (bo agent ma 60 lat, a wiek bazowy to 30).

Model Gompertza:
\[
\lambda_{\text{CVD}}(a) = \lambda_{0,\text{CVD}} \cdot \exp(0.06 \cdot (a - 30))
\]

Dla $a = 60$: $\exp(0.06 \cdot 30) = e^{1.8} \approx 6.05$. Znaczy to, że hazard CVD u 60-latka jest 6.05$\times$ wyższy niż u 30-latka (oba bez RF).

\textbf{Czas podwojenia hazardu:}
\[
T_{\text{podwojenia}} = \frac{\ln 2}{\gamma_{\text{age}}} = \frac{0.693}{0.06} \approx 11.55 \text{ lat}
\]

\subsubsection{Współczynniki $\beta = \ln(\text{HR})$ z macierzy \texttt{HAZARD\_BETA}}

\begin{table}[h]
\centering
\small
\begin{tabular}{lrrl}
\toprule
\textbf{Risk Factor} & \textbf{HR} & \textbf{$\beta = \ln(\text{HR})$} & \textbf{Źródło} \\
\midrule
smoking            & 2.5 & 0.916 & ESC, Framingham \\
high\_cholesterol  & 2.0 & 0.693 & ATP III, Framingham \\
hypertension       & 2.2 & 0.788 & Wytyczne ESC \\
\bottomrule
\end{tabular}
\caption{HR i $\beta$ użyte w przykładzie.}
\end{table}

Suma wykładników: $0.916 + 0.693 + 0.788 = 2.397$, więc $\exp(2.397) \approx 10.99$.

\subsection{Pełne obliczenie krok po kroku}

\begin{align*}
\Delta h &= \underbrace{6.0 \cdot 10^{-5}}_{\lambda_0} \cdot \underbrace{6.05}_{\text{wiek}} \cdot \underbrace{10.99}_{\text{RF}} \\
         &\approx 3.99 \cdot 10^{-3} \text{ /miesiąc}
\end{align*}

\textbf{Porównanie z baseline} (60-latek bez RF):
\[
\Delta h_{\text{baseline}} = 6.0 \cdot 10^{-5} \cdot 6.05 = 3.63 \cdot 10^{-4}
\]
\[
\text{Mnożnik} = \frac{\Delta h}{\Delta h_{\text{baseline}}} = \frac{3.99 \cdot 10^{-3}}{3.63 \cdot 10^{-4}} \approx 11{,}0\times
\]

Z aproksymacji Poissona:
\[
P(\text{onset CVD w tym miesiącu}) = 1 - e^{-\Delta h} \approx 0.40\,\%
\]
W skali roku: $1 - (1 - 0.004)^{12} \approx 4.7\,\%$.

\subsection{Dlaczego wiek bazowy to 30 lat?}

Wybór 30 lat jako wieku referencyjnego $a_{\text{ref}}$ w modelu Gompertza:
\[
\lambda(a) = \lambda_0 \cdot \exp(\gamma_{\text{age}} \cdot (a - a_{\text{ref}}))
\]

\begin{enumerate}[leftmargin=*, itemsep=4pt]
  \item \textbf{Konwencja epidemiologiczna.} W aktuarialnej i epidemiologicznej praktyce 30 lat jest standardowym wiekiem referencyjnym dla chronic diseases. Większość badań kohortowych (Framingham, MRFIT) zaczyna pomiary między 25 a 35 r.ż.

  \item \textbf{Przed 30 r.ż.\ chronic disease hazard jest pomijalny.} Dominują przyczyny niezwiązane z RF: wypadki, samobójstwa, zakażenia. Dopiero po 30 r.ż.\ Gompertz exponential growth zaczyna dominować.

  \item \textbf{Numerycznie wygodny.} $a - 30$ zachowuje stosunkowo niewielkie wartości wykładnika dla typowych wieków agentów (do $a = 100$, $a - 30 = 70$, $\exp(0.06 \cdot 70) = e^{4.2} \approx 67$ --- jeszcze poniżej numerycznego overflow).

  \item \textbf{Symetria interpretacyjna.} Średnia długość życia człowieka to $\sim$78--82 lata; wiek 30 to $\sim$1/3 życia, dobrze reprezentuje "młodszy dorosły".

  \item \textbf{Brak hazardu u dzieci.} W kodzie dodatkowo:
  \[
  \text{effective\_age} = \max(a,\ 18)
  \]
  Dzięki temu agenci $<$ 18 r.ż.\ nie kumulują $H_{\mathrm{cum}}$ --- biologicznie poprawne (chronic diseases nie rozwijają się u dzieci poza specyficznymi syndromami).
\end{enumerate}

\subsection{Co by było, gdyby wiek bazowy był inny?}

\begin{warnbox}
Wartość $a_{\text{ref}} = 30$ jest \emph{arbitralna} matematycznie --- ma sens fizyczny, ale model jest \emph{niezmienniczy} na zmianę punktu referencyjnego, o ile odpowiednio przeskalujemy $\lambda_0$. To znaczy: jeśli zmienimy $a_{\text{ref}} = 20$, ale jednocześnie pomnożymy $\lambda_0$ przez $\exp(\gamma_{\text{age}} \cdot 10)$, dostaniemy identyczne hazardy. Konwencja 30 jest po prostu czytelna i powszechna.
\end{warnbox}

\newpage
% ============================================================
\section{CBR walidacji: 8.30 vs 8.50 promila --- skąd różnica?}
% ============================================================

\begin{qbox}
W tabeli walidacji jest napisane CBR\,=\,8.30\,\promile{} dla modelu, a u Polski 8.50\,\promile{}. Skąd ta drobna różnica? Czy to błąd?
\end{qbox}

\subsection{Wartości referencyjne i modelowe}

\begin{table}[h]
\centering
\small
\begin{tabular}{lrrl}
\toprule
\textbf{Wskaźnik} & \textbf{Polska 2021 (GUS)} & \textbf{Model (FM=MM=1.0)} & \textbf{Zgodność} \\
\midrule
CBR        & 8.50\promile  & \textbf{8.30\promile}  & 97.6\,\% (drobne niedoszacowanie) \\
CDR        & 13.50\promile & \textbf{15.60\promile} & 115\,\% (zawyżone, patrz \ref{ssec:cdr}) \\
TFR        & 1.26          & \textbf{1.28}          & 101.6\,\% (zgodność wzorcowa) \\
\bottomrule
\end{tabular}
\caption{Porównanie modelu z danymi GUS Polska 2021.}
\end{table}

\subsection{Dlaczego CBR jest 8.30 zamiast 8.50?}

Kilka przyczyn drobnej różnicy (0.2\promile, czyli 2.4\,\% wartości):

\begin{enumerate}[leftmargin=*, itemsep=4pt]
  \item \textbf{Zaokrąglenia w tabeli ASFR.} Wartości ASFR w naszym kodzie są zaokrąglone do trzech miejsc po przecinku:
  \begin{lstlisting}
DEFAULT_FERTILITY_TABLE = {
    15: 0.011,  20: 0.041,  25: 0.081,
    30: 0.082,  35: 0.034,  40: 0.007,
    45: 0.0003,
}
\end{lstlisting}
  GUS publikuje wartości z większą precyzją; te zaokrąglenia kumulują się do $\sim$1\,\% niedoszacowania TFR (1.28 vs 1.26 --- już sam ten brak skali pokrywa większość różnicy).

  \item \textbf{Dyskretyzacja miesięczna.} Polski roczny ASFR jest dzielony na 12, ale empirycznie urodzenia nie rozkładają się równomiernie (są lekko sezonowe). To wprowadza drobne odchylenie $\sim$0.1\promile.

  \item \textbf{Synteza populacji.} Populacja syntetyczna 50\,000 agentów jest \emph{losowana} z rozkładu wiekowego GUS. Każde uruchomienie ma inny rozkład, średnio bliski celowi, ale z wariancją.

  \item \textbf{Korekta przez choroby (\texttt{disease\_reduction}).} W modelu:
  \[
  P_{\text{birth}} = \frac{\text{ASFR}}{12} \cdot \text{FM} \cdot \underbrace{\max(1 - 0.02 \cdot n - 0.04 \cdot \text{dis},\ 0.7)}_{\geq\, 0.7}
  \]
  Część kobiet ma aktywne choroby (CVD lub Lung Cancer), co lekko redukuje płodność. To efekt nieobecny w danych GUS, które dotyczą "uśrednionej" Polski.
\end{enumerate}

\begin{infobox}
\textbf{Wniosek.} Różnica 0.2\promile{} (2.4\,\%) mieści się w tolerancji modelu. CBR\,=\,8.30 vs GUS 8.50 --- to wartości \emph{w pełni akceptowalne} dla symulacji demograficznej. Dokładność lepsza niż 5\,\% w demografii to bardzo dobry wynik.
\end{infobox}

\subsection{Dlaczego CDR jest 15.6 zamiast 13.5? (poważniejsza rozbieżność)}
\label{ssec:cdr}

To większa rozbieżność (15\,\%) i ma inną przyczynę:

\begin{warnbox}
\textbf{Powód:} Tabele $q_x$ z GUS dotyczą \emph{całej} polskiej populacji --- włącznie z osobami przewlekle chorymi, których ryzyko zgonu jest już "wbudowane" w średnie stopy.

W naszym modelu populacja syntetyczna na starcie ma 35\,\% prevalencji CVD, ale nie ma populacyjnego "rezerwuaru" osób z wieloma chorobami współistniejącymi. Tabela $q_x$ jest aplikowana do agentów, którzy są przeciętnie \emph{zdrowsi} niż polska populacja referencyjna --- a mimo to dostają pełną $q_x$, plus dodatkowe \texttt{cox\_mult} dla aktywnych chorób. To podwójnie zwiększa mortality.

\textbf{Akceptowalność:} odchylenie jest świadome i udokumentowane. Ważne, że punkt stabilności (FM, MM) pozostaje wewnątrz badanego zakresu siatki.
\end{warnbox}

\subsection{Dlaczego TFR jest niemal idealne (1.28 vs 1.26)?}

TFR = $\sum \text{ASFR} \cdot \Delta\text{age}$ --- to czysta suma wartości tabelarycznych, bez wpływu struktury wiekowej populacji ani innych efektów dynamicznych. Mała różnica (1.6\,\%) wynika tylko z zaokrągleń tabeli ASFR.

\subsection{Inżynierski wniosek}

W demograficznym ABM dokładność lepsza niż 5\,\% dla CBR i niż 20\,\% dla CDR jest \emph{bardzo dobra}. Model:
\begin{itemize}
  \item Replikuje TFR z dokładnością $<$2\,\% (idealnie).
  \item Replikuje CBR z dokładnością $<$3\,\% (bardzo dobrze).
  \item Replikuje CDR z dokładnością 15\,\% (akceptowalnie --- z udokumentowanym powodem).
\end{itemize}

\newpage
% ============================================================
\section{Co oznacza \texttt{Dict[int, Citizen]}? Co to jest \texttt{Zone}?}
% ============================================================

\begin{qbox}
W pseudokodzie i diagramie architektury pojawiają się \texttt{Dict[int, Citizen]} oraz klasa \texttt{Zone}. Co to dokładnie znaczy?
\end{qbox}

\subsection{\texttt{Dict[int, Citizen]} --- typowa adnotacja Pythona}

Jest to \textbf{adnotacja typu} (type hint) z modułu \texttt{typing}. Oznacza:

\begin{formulabox}
\texttt{Dict[int, Citizen]} = \emph{słownik}, w którym:
\begin{itemize}
  \item \textbf{klucz} jest typu \texttt{int} (liczba całkowita)
  \item \textbf{wartość} jest typu \texttt{Citizen} (obiekt klasy Citizen)
\end{itemize}
\end{formulabox}

W kontekście modelu, to jest \emph{rejestr populacji}:

\begin{lstlisting}
class SimulationEngine:
    def __init__(self, ...):
        self.citizens: Dict[int, Citizen] = {}
\end{lstlisting}

Przykład zawartości:
\begin{lstlisting}
self.citizens = {
    1:    Citizen(id=1,    sex="male",   age_months=384, alive=True,  ...),
    2:    Citizen(id=2,    sex="female", age_months=252, alive=True,  ...),
    3:    Citizen(id=3,    sex="female", age_months=816, alive=False, ...),
    ...
    50000: Citizen(id=50000, ...),
}
\end{lstlisting}

\subsubsection{Klucz = ID agenta}

Każdy obywatel ma unikalne ID generowane przez zmienną klasową \texttt{Citizen.\_id\_counter}, automatycznie inkrementowaną przy tworzeniu obiektu. ID jest niezmienne przez całe życie agenta i nie powtarza się nawet po jego śmierci.

\subsubsection{Wartość = obiekt Citizen}

Każdy obiekt klasy \texttt{Citizen} zawiera kompletny stan agenta:

\begin{lstlisting}
class Citizen:
    id: int                              # unikalny identyfikator
    sex: str                             # "male" lub "female"
    age_months: int                      # wiek w miesiacach (nie latach!)
    alive: bool                          # True/False
    household_id: int                    # przynaleznosc do gospodarstwa
    zone_id: int                         # strefa geograficzna (1-4)
    diseases: Dict[str, int]             # {"CVD": 0, "Lung Cancer": 1}
    risk_factors: Dict[str, int]         # 7 RF: {"smoking": 0, ...}
    cumulative_hazard: Dict[str, float]  # pamiec Coxa
    disability_score: float              # suma wag chorob
\end{lstlisting}

\subsubsection{Dlaczego Dict, a nie List?}

\begin{table}[h]
\centering
\small
\begin{tabular}{lll}
\toprule
\textbf{Kryterium} & \texttt{Dict[int, Citizen]} & \texttt{List[Citizen]} \\
\midrule
Lookup po ID                & $O(1)$               & $O(n)$ \\
Usunięcie martwych          & nie trzeba (flaga \texttt{alive}) & drogie przesunięcia \\
Iteracja po żywych          & \texttt{if c.alive}  & analogicznie \\
Nowe narodziny              & \texttt{citizens[id] = newborn} & \texttt{append()} \\
Stabilność referencji       & stałe ID & indeksy się przesuwają \\
\bottomrule
\end{tabular}
\caption{Porównanie struktur danych dla populacji.}
\end{table}

\textbf{Kluczowe:} martwi agenci \emph{nie są usuwani} ze słownika. Mają flagę \texttt{alive = False}. Dzięki temu:
\begin{itemize}
  \item gospodarstwo może sprawdzić czy członek żyje bez \texttt{KeyError},
  \item statystyki retrospektywne są możliwe (np.\ ile osób zmarło w wieku 60--70 lat),
  \item unika się kosztownej przebudowy słownika w trakcie symulacji.
\end{itemize}

\subsection{Klasa \texttt{Zone} --- strefy geograficzne}

\texttt{Zone} reprezentuje \textbf{strefę geograficzną} z parametrami środowiskowymi. Model ma 4 strefy, każda z własnymi parametrami.

\subsubsection{Definicja klasy}

\begin{lstlisting}
class Zone:
    id: int                  # identyfikator (1-4)
    air_quality: float       # 0.0 (zlej) - 1.0 (doskonalej)
    healthcare_access: float # 0.0 - 1.0, dostepnosc opieki
    density: int             # gestosc zaludnienia (osob/km^2)
\end{lstlisting}

\subsubsection{Cztery strefy w modelu}

\begin{table}[h]
\centering
\small
\begin{tabular}{rrrr}
\toprule
\textbf{Strefa} & \textbf{Jakość powietrza} & \textbf{Dostęp do opieki} & \textbf{Gęstość (osób/km$^2$)} \\
\midrule
1 & 0.75 (wysoka)  & 0.90 (dobry)           & 4000 \\
2 & 0.65 (średnia) & 0.85                   & 6000 \\
3 & 0.70           & 0.80                   & 5000 \\
4 & 0.55 (niska)   & 0.75 (ograniczony)     & 8000 \\
\bottomrule
\end{tabular}
\caption{Parametry 4 stref modelu.}
\end{table}

Każdy agent ma atrybut \texttt{zone\_id} wskazujący na jedną z tych stref. Gospodarstwa domowe (\texttt{Household}) również należą do stref.

\subsubsection{Reprezentacja w \texttt{SimulationEngine}}

\begin{lstlisting}
class SimulationEngine:
    def __init__(self, ...):
        self.citizens:   Dict[int, Citizen]   = {}
        self.households: Dict[int, Household] = {}
        self.zones:      Dict[int, Zone]      = {
            1: Zone(id=1, air_quality=0.75, ...),
            2: Zone(id=2, air_quality=0.65, ...),
            3: Zone(id=3, air_quality=0.70, ...),
            4: Zone(id=4, air_quality=0.55, ...),
        }
\end{lstlisting}

\subsubsection{Wpływ stref na model (obecnie minimalny)}

\begin{warnbox}
\textbf{Status w obecnej wersji:} parametry środowiskowe stref są \emph{przechowywane}, ale ich wpływ na śmiertelność i choroby jest \emph{minimalny}. Głównym motorem demograficznym są wiekowe tabele GUS (zgodnie z dokumentacją projektu).

\textbf{Potencjalne rozszerzenie:} kalibrowane wagi mogłyby aktywować różnice między strefami:
\begin{itemize}
  \item Strefa 4 (wysoka gęstość, niska jakość powietrza) --- wyższa baza dla \texttt{smoking}, \texttt{respiratory disease}.
  \item Strefa 1 (przedmieścia, niska gęstość) --- niższa baza dla \texttt{air-pollution-related} ryzyka.
\end{itemize}
W obecnej wersji projekt skupia się na demografii bazowej + 7 RF (model Coxa) jako głównym driverze.
\end{warnbox}

\subsection{Pełna architektura w jednym obrazie}

\begin{lstlisting}
SimulationEngine
|-- citizens:   Dict[int, Citizen]     # 50 000 obywateli, klucz = id
|-- households: Dict[int, Household]   # rodziny, klucz = id
|-- zones:      Dict[int, Zone]        # 4 strefy
|-- disease_model: DiseaseModel        # parametry Coxa
|-- yearly_stats:  Dict[int, Dict]     # statystyki per rok

Citizen
|-- id, sex, age_months, alive
|-- household_id     -> klucz do households
|-- zone_id          -> klucz do zones
|-- diseases:        Dict[str, int]
|-- risk_factors:    Dict[str, int]    # 7 RF
|-- cumulative_hazard: Dict[str, float]  # pamiec Coxa

Household
|-- id, zone_id
|-- members: List[int]   # lista ID agentow (nie obiektow!)

Zone
|-- id, air_quality, healthcare_access, density
\end{lstlisting}

\textbf{Ważny szczegół:} \texttt{Household.members} przechowuje \emph{listę ID} obywateli, nie referencji do obiektów. To zapobiega cyklicznym zależnościom i ułatwia serializację (np.\ zapis do JSON).

\vfill
\begin{center}
\color{muted}\small
\rule{0.4\textwidth}{0.4pt}\\[8pt]
Notatki uzupełniające do projektu \emph{Urban Health ABM}\\
ICM UW \ \textbar\ \ 2026
\end{center}

\end{document}
